\begin{center}\begin{tikzpicture}[font=\small,>=stealth,line join=round]
\draw[gray!40,rounded corners=3pt] (0,0) rectangle (14.2,8.6);
\node at (7.1,8.0) {例0与快练：两个最小位置相同，可以同时取到};
\begin{scope}[shift={(4.2,4.25)},x=0.8235294117647058cm,y=0.8235294117647058cm]
\draw[->,gray] (-1,0)--(3,0) node[right]{$x$};
\draw[blue,very thick] (0,-2)--(0,3);
\draw[orange!85!black,dashed,thick] (0,-2)--(2,0)--(0,3);
\draw[dotted,gray] (0,0)--(0,-2) (0,0)--(0,3);
\fill (0,-2) circle(2pt) (0,3) circle(2pt);
\fill[blue] (0.0,0) circle(2.3pt);\fill[orange!85!black] (2,0) circle(2pt);
\node[below right=5pt] at (0,-2) {$A$};\node[above right=5pt] at (0,3) {$B$};
\node[below right=8pt] at (0.0,0) {$P_0$};\node[above right=6pt] at (2,0) {$P$};
\end{scope}
\node[align=left,anchor=west] at (8.2,4.9) {$A=(0,-2),\ B=(0,3)$\\[12pt]$P_0=(0,0)$\\[12pt]最短路程 $2+3=5$\\[14pt]一般 $m=n=c$ 时\\[9pt]$x_0=c,\ f_{\min}=a+b$};
\node[text width=12.8cm,align=center] at (7.1,.65) {横纵坐标等比例；蓝线为最短直路，橙色虚线为经其他位置 $P$ 的路。};
\end{tikzpicture}\end{center}
